% !TEX root = ../thesis.tex

% 中英文摘要和关键字

\begin{abstract}
  % 原中文摘要（润色前）：
  % 大尺度结构观测中的宇宙学参数推断依赖大量模拟样本估计协方差矩阵，并需要同时刻画暗物质密度场和速度场。传统 N 体模拟精度高但计算代价昂贵，基于拉格朗日微扰论的快速方法计算效率高，却难以恢复非线性小尺度结构和速度场细节。针对这一问题，本文提出一种数据驱动的宇宙学快速模拟流程，以 Zel'dovich 近似（ZA）提供位移场和速度场的大尺度基准，再使用两个三维生成对抗残差网络分别学习位移残差和速度残差。位移网络以 ZA 位移场为条件输入；速度网络以 ZA 速度场和位移网络推理得到的最终位移场为联合条件输入，从而利用已恢复的非线性结构信息补全速度场。网络采用三维 U 形生成器和多尺度判别器，并结合分块推理、高通滤波和旋转平均策略，使方法能够自然地扩展到同分辨率、更大体积的模拟中。
  %
  % 本文使用 Quijote N 体模拟验证该流程。结果表明，与 ZA、2LPT 和 ALPT 等微扰论方法相比，新流程在暗物质密度场的切片形态、概率密度函数、功率谱和双谱上均更接近 N 体参考结果，能够显著恢复小尺度非线性功率和高阶关联性。功率谱协方差矩阵的整体结构也与 N 体结果较为一致，但小尺度对角方差仍存在一定低估，且当前估计受样本数量限制。速度场方面，新流程在速度模切片、速度散度功率谱和各速度分量功率谱上相较 ZA 有明显改进，但对旋度和维里化区域中的小尺度随机速度仍恢复不足。总体而言，本文方法在保持大尺度结构一致性的同时改善了密度场和速度场的小尺度统计性质，为面向多探针宇宙学分析的快速模拟提供了可行方案。

  大尺度结构观测中的宇宙学参数推断通常需要大量模拟样本来估计协方差矩阵，并要求模拟能够同时给出可靠的暗物质密度场和速度场。传统 N 体模拟具有较高精度，但计算代价高昂；基于拉格朗日微扰论的快速模拟方法效率较高，却难以准确刻画非线性小尺度结构及速度场细节。针对这一问题，本文提出一种数据驱动的宇宙学快速模拟流程：首先以 Zel'dovich 近似（ZA）生成位移场和速度场的大尺度基准，再分别训练两个三维生成对抗网络学习位移残差和速度残差。其中，位移网络以 ZA 位移场为条件输入；速度网络以 ZA 速度场和位移网络推理得到的最终位移场为联合条件输入，从而利用已恢复的非线性结构信息进一步修正速度场。网络采用三维 U 形生成器与多尺度判别器，并在推理阶段结合分块推理、高通滤波和旋转平均等策略，使该流程在保持快速生成优势的同时，能够稳定扩展至同分辨率下更大体积的模拟任务。

  本文基于 Quijote N 体模拟对所提出流程进行了验证。结果表明，与 ZA、2LPT 和 ALPT 等微扰论方法相比，本文方法生成的暗物质密度场在切片形态、概率密度函数、功率谱、双谱和功率谱协方差矩阵等统计量上均更接近 N 体参考结果，能够有效恢复小尺度非线性功率和高阶关联。在速度场方面，本文方法相较 ZA 在速度模切片、速度散度功率谱以及各速度分量功率谱上均有明显改进，对旋度以及维里化区域中的小尺度随机速度也有一定恢复。总体而言，本文方法在保持大尺度结构一致性的同时，改善了密度场和速度场的小尺度统计性质，为面向多探针宇宙学分析的快速模拟提供了一种可行方案。

  % 关键词用“英文逗号”分隔，输出时会自动处理为正确的分隔符
  \thusetup{
    keywords = {宇宙学快速模拟, N 体模拟, 生成对抗网络, 暗物质密度场, 速度场},
  }
\end{abstract}

\begin{abstract*}
  Cosmological parameter inference from large-scale structure observations generally requires a large number of simulation realizations for covariance-matrix estimation, and these simulations should provide reliable dark matter density and velocity fields at the same time. Full N-body simulations are accurate but computationally expensive. Fast simulation methods based on Lagrangian perturbation theory are much more efficient, but they cannot accurately describe nonlinear small-scale structures or detailed velocity-field features. To address this problem, this thesis proposes a data-driven fast cosmological simulation pipeline. The pipeline first uses the Zel'dovich approximation (ZA) to generate large-scale reference fields for displacement and velocity, and then trains two three-dimensional generative adversarial networks to learn the corresponding displacement and velocity residuals. The displacement network takes the ZA displacement field as its conditional input. The velocity network is conditioned jointly on the ZA velocity field and the final displacement field inferred by the displacement network, so that the recovered nonlinear structural information can be used to further correct the velocity field. The networks adopt three-dimensional U-shaped generators and multi-scale discriminators. At the inference stage, tiled inference, high-pass filtering, and rotation averaging are combined, enabling the pipeline to scale stably to larger simulation volumes at the same resolution while preserving the advantage of fast generation.

  The proposed pipeline is validated using Quijote N-body simulations. The results show that, compared with perturbative methods such as ZA, 2LPT, and ALPT, the dark matter density fields generated by the proposed method are closer to the N-body reference in slice morphology, probability density functions, power spectra, bispectra, and power-spectrum covariance matrices. The method therefore effectively recovers nonlinear small-scale power and higher-order correlations. For the velocity field, the proposed method achieves clear improvements over ZA in velocity-magnitude slices, velocity-divergence power spectra, and power spectra of individual velocity components. It also provides partial recovery of vorticity and small-scale stochastic velocities in virialized regions. Overall, the proposed method improves the small-scale statistical properties of both density and velocity fields while maintaining consistency on large scales, providing a feasible fast-simulation approach for multi-probe cosmological analyses.

  % Use comma as separator when inputting
  \thusetup{
    keywords* = {fast cosmological simulation, N-body simulation, generative adversarial network, dark matter density field, velocity field},
  }
\end{abstract*}
